\chapter{Repaso y evaluación A1}
\section{Objetivos}
\begin{itemize}[leftmargin=*]
	\item Consolidar los contenidos de los capítulos 1--19.
	\item Preparar una autoevaluación tipo checklist A1.
	\item Practicar una mini prueba integrada (comprensión, gramática y producción).
\end{itemize}

\section{Componentes del repaso}
\begin{itemize}[leftmargin=*]
	\item \textbf{Gramática}: ejercicios mixtos de V2, negación, modales y orden con infinitivo final.
	\item \textbf{Léxico}: campos semánticos principales (familia, comida, ciudad, rutina, tiempo).
	\item \textbf{Pronunciación}: repaso de sonidos difíciles y ritmo de frase.
\end{itemize}

\section{Ruta del capítulo}
Siete secciones para repasar gramática, vocabulario y destrezas con una simulación final.

\section{Repaso de gramática A1}
\textbf{Temas clave.}
\begin{itemize}[leftmargin=*]
  \item Orden V2 en afirmativas y preguntas sí/no. 
  \item Negación con \textit{niet/geen}. 
  \item Modales + infinitivo final. 
  \item Presente de regulares e irregulares básicos. 
  \item Artículos \textit{de/het/een} y plural. 
\end{itemize}

\textbf{Mini práctica}
\begin{itemize}[leftmargin=*]
  \item 10 frases mixtas para corregir (V2, negación, modales).
  \item Pequeño dictado de 5 frases centrado en orden de palabras.
\end{itemize}

\section{Repaso de vocabulario A1}
\textbf{Campos semánticos.}
\begin{itemize}[leftmargin=*]
  \item Familia, comida/bebida, ciudad/transporte, rutina diaria, tiempo/horarios. 
  \item Colores, adjetivos básicos, números 0--100. 
\end{itemize}

\textbf{Actividades}
\begin{itemize}[leftmargin=*]
  \item Listas rápidas de 8 palabras por campo y uso en frase. 
  \item Juego de definiciones: describe en neerlandés, otro adivina la palabra. 
\end{itemize}

\section{Comprensión auditiva básica}
\textbf{Práctica sugerida.}
\begin{itemize}[leftmargin=*]
  \item Escucha 2 audios cortos (30--45s) sobre compras y transporte; responde 3 preguntas cada uno.
  \item Usa preguntas de opción múltiple o verdadero/falso.
  \item Repite en voz alta frases clave para reforzar ritmo.
\end{itemize}

\textbf{Mini práctica}
\begin{itemize}[leftmargin=*]
  \item Transcribe 4 frases del audio y compáralas con la solución.
  \item Marca palabras nuevas y pásalas a flashcards.
\end{itemize}

\section{Producción oral guiada}
\textbf{Tareas cortas.}
\begin{itemize}[leftmargin=*]
  \item Monólogo de 2 minutos: presentación + rutina + plan de fin de semana.
  \item Diálogo de pedido en café (4 líneas por persona).
  \item Preguntas rápidas (10): responde en 1 frase con V2 correcto.
\end{itemize}

\textbf{Autoevaluación}
\begin{itemize}[leftmargin=*]
  \item Revisa pronunciación de \textit{g/ch, ui, ij/ei}.
  \item Comprueba ritmo y claridad; graba y escucha.
\end{itemize}

\section{Lectura A1}
\textbf{Textos modelo.}
\begin{itemize}[leftmargin=*]
  \item Anuncio corto (transporte o tienda) de unas~80 palabras.
  \item Email simple entre amigos sobre planes de fin de semana.
\end{itemize}

\textbf{Tareas}
\begin{itemize}[leftmargin=*]
  \item Subraya verbos y preposiciones de tiempo/lugar.
  \item Responde 4 preguntas de comprensión por texto.
\end{itemize}

\section{Escritura básica}
\textbf{Tareas cortas.}
\begin{itemize}[leftmargin=*]
  \item Correo de 80--100 palabras: presentarte, decir dónde vives, qué haces y tu plan del sábado.
  \item Nota rápida para un amigo (40 palabras) cancelando o moviendo una cita.
\end{itemize}

\textbf{Checklist}
\begin{itemize}[leftmargin=*]
  \item Orden V2, artículos correctos, negación sencilla, preposiciones de tiempo/lugar.
  \item Releer y simplificar oraciones largas.
\end{itemize}

\section{Simulación de examen A1}
\textbf{Estructura sugerida (45--60 min).}
\begin{itemize}[leftmargin=*]
  \item Comprensión auditiva: 2 audios cortos, 6 preguntas en total.
  \item Lectura: 1 texto breve con 6 preguntas (Opción múltiple o VF).
  \item Gramática: 10 ítems mixtos (V2, negación, modales, artículos).
  \item Escritura: correo de 80 palabras siguiendo un prompt.
  \item Oral: monólogo de 2 minutos o diálogo de 8 turnos (grabar).
\end{itemize}

\textbf{Rúbrica rápida}
\begin{itemize}[leftmargin=*]
  \item Claridad (pronunciación), precisión (gramática/artículos), adecuación (vocabulario A1), coherencia (orden lógico).
\end{itemize}


\section{Prueba modelo}
\begin{itemize}[leftmargin=*]
	\item Comprensión auditiva: 2 audios cortos con preguntas de opción múltiple.
	\item Lectura: texto breve con 6 preguntas.
	\item Producción escrita: correo simple (80--100 palabras) presentándote y contando tu día.
	\item Producción oral: monólogo de 2 minutos o diálogo grabado.
\end{itemize}

\section{Siguiente paso}
\begin{itemize}[leftmargin=*]
	\item Identifica vacíos y diseña un plan de 4 semanas para pasar a nivel A2 o reforzar A1.
\end{itemize}
